%% Paper A -- Engineering Applications of Artificial Intelligence
%% Cooperative and non-cooperative game-theoretic scheduling of production machines and
%% battery-constrained autonomous vehicles under time-of-use tariffs.
%%
%% Every numeral in a results sentence is written as a jnum macro and resolves against
%% numbers.tex, generated by experiments/run_all.py from recorded runs.  A literal numeral
%% in a results sentence is a defect (referee check R-5).

\documentclass[preprint, 3p, review, 12pt]{elsarticle}

\usepackage{amssymb}
\usepackage{amsmath}
\usepackage{amsthm}
\usepackage{adjustbox}
\usepackage{algpseudocode}
\usepackage{algorithm}
\usepackage{caption, comment}
\usepackage{subcaption, multirow, blkarray, tabularx}
\usepackage{graphicx}
\usepackage{lscape, booktabs}
\usepackage{url}
\usepackage{xcolor}
\usepackage{textcomp}
\providecommand{\euro}{\texteuro}

\newtheorem{theorem}{Theorem}
\newtheorem{proposition}[theorem]{Proposition}
\newtheorem{corollary}[theorem]{Corollary}
\newtheorem{observation}[theorem]{Observation}
\theoremstyle{definition}
\newtheorem{definition}[theorem]{Definition}
\newtheorem{assumption}[theorem]{Assumption}

% Generated by jsspt_tou.analysis.registry -- do not edit by hand.
\makeatletter
\newcommand{\jnum}[1]{\@ifundefined{jnum@#1}{\textbf{??#1??}}{\csname jnum@#1\endcsname}}
\makeatother
\expandafter\def\csname jnum@bench.instances\endcsname{40}
\expandafter\def\csname jnum@bench.omegas\endcsname{3}
\expandafter\def\csname jnum@e1.droracle.cmax.mean\endcsname{129.8}
\expandafter\def\csname jnum@e1.droracle.deadline.rate\endcsname{100.0}
\expandafter\def\csname jnum@e1.droracle.ecost.mean\endcsname{1.389}
\expandafter\def\csname jnum@e1.droracle.fallback.rate\endcsname{0.0}
\expandafter\def\csname jnum@e1.droracle.phi.mean\endcsname{0.2564}
\expandafter\def\csname jnum@e1.droracle.phi.std\endcsname{0.0588}
\expandafter\def\csname jnum@e1.droracle.seconds.mean\endcsname{0.003}
\expandafter\def\csname jnum@e1.drspt.cmax.mean\endcsname{163.9}
\expandafter\def\csname jnum@e1.drspt.deadline.rate\endcsname{100.0}
\expandafter\def\csname jnum@e1.drspt.ecost.mean\endcsname{1.485}
\expandafter\def\csname jnum@e1.drspt.fallback.rate\endcsname{0.0}
\expandafter\def\csname jnum@e1.drspt.phi.mean\endcsname{0.3974}
\expandafter\def\csname jnum@e1.drspt.phi.std\endcsname{0.1194}
\expandafter\def\csname jnum@e1.drspt.seconds.mean\endcsname{0.003}
\expandafter\def\csname jnum@e1.drtou.cmax.mean\endcsname{251.5}
\expandafter\def\csname jnum@e1.drtou.deadline.rate\endcsname{27.5}
\expandafter\def\csname jnum@e1.drtou.ecost.mean\endcsname{0.853}
\expandafter\def\csname jnum@e1.drtou.fallback.rate\endcsname{72.5}
\expandafter\def\csname jnum@e1.drtou.phi.mean\endcsname{2.5175}
\expandafter\def\csname jnum@e1.drtou.phi.std\endcsname{1.8313}
\expandafter\def\csname jnum@e1.drtou.seconds.mean\endcsname{0.003}
\expandafter\def\csname jnum@e1.friedman.cd\endcsname{1.425}
\expandafter\def\csname jnum@e1.friedman.p\endcsname{0.000000}
\expandafter\def\csname jnum@e1.friedman.stat\endcsname{170.78}
\expandafter\def\csname jnum@e1.m1a.cmax.mean\endcsname{124.6}
\expandafter\def\csname jnum@e1.m1a.deadline.rate\endcsname{100.0}
\expandafter\def\csname jnum@e1.m1a.ecost.mean\endcsname{1.348}
\expandafter\def\csname jnum@e1.m1a.epsne.max\endcsname{0.000000000}
\expandafter\def\csname jnum@e1.m1a.fallback.rate\endcsname{0.0}
\expandafter\def\csname jnum@e1.m1a.gain.mean\endcsname{22.1}
\expandafter\def\csname jnum@e1.m1a.gain.std\endcsname{15.0}
\expandafter\def\csname jnum@e1.m1a.iterations.mean\endcsname{2.06}
\expandafter\def\csname jnum@e1.m1a.phi.mean\endcsname{0.2344}
\expandafter\def\csname jnum@e1.m1a.phi.std\endcsname{0.0610}
\expandafter\def\csname jnum@e1.m1a.seconds.mean\endcsname{0.612}
\expandafter\def\csname jnum@e1.m1a.stages.mean\endcsname{27.5}
\expandafter\def\csname jnum@e1.m1a.winrate.vs.oracle\endcsname{65.0}
\expandafter\def\csname jnum@e1.m1b.cmax.mean\endcsname{149.8}
\expandafter\def\csname jnum@e1.m1b.deadline.rate\endcsname{100.0}
\expandafter\def\csname jnum@e1.m1b.ecost.mean\endcsname{1.361}
\expandafter\def\csname jnum@e1.m1b.fallback.rate\endcsname{0.0}
\expandafter\def\csname jnum@e1.m1b.phi.mean\endcsname{0.3314}
\expandafter\def\csname jnum@e1.m1b.phi.std\endcsname{0.0999}
\expandafter\def\csname jnum@e1.m1b.seconds.mean\endcsname{0.954}
\expandafter\def\csname jnum@e1.m2grand.cmax.mean\endcsname{123.0}
\expandafter\def\csname jnum@e1.m2grand.deadline.rate\endcsname{100.0}
\expandafter\def\csname jnum@e1.m2grand.ecost.mean\endcsname{1.324}
\expandafter\def\csname jnum@e1.m2grand.fallback.rate\endcsname{0.0}
\expandafter\def\csname jnum@e1.m2grand.gain.mean\endcsname{24.2}
\expandafter\def\csname jnum@e1.m2grand.gain.std\endcsname{17.0}
\expandafter\def\csname jnum@e1.m2grand.phi.mean\endcsname{0.2256}
\expandafter\def\csname jnum@e1.m2grand.phi.std\endcsname{0.0513}
\expandafter\def\csname jnum@e1.m2grand.seconds.mean\endcsname{0.263}
\expandafter\def\csname jnum@e1.m2grand.winrate.vs.oracle\endcsname{75.0}
\expandafter\def\csname jnum@e1.pi0.cmax.mean\endcsname{134.6}
\expandafter\def\csname jnum@e1.pi0.deadline.rate\endcsname{100.0}
\expandafter\def\csname jnum@e1.pi0.ecost.mean\endcsname{1.594}
\expandafter\def\csname jnum@e1.pi0.fallback.rate\endcsname{0.0}
\expandafter\def\csname jnum@e1.pi0.phi.mean\endcsname{0.3077}
\expandafter\def\csname jnum@e1.pi0.phi.std\endcsname{0.0768}
\expandafter\def\csname jnum@e1.pi0.seconds.mean\endcsname{0.001}
\expandafter\def\csname jnum@e1.rank.droracle\endcsname{2.89}
\expandafter\def\csname jnum@e1.rank.drspt\endcsname{5.31}
\expandafter\def\csname jnum@e1.rank.drtou\endcsname{6.92}
\expandafter\def\csname jnum@e1.rank.m1a\endcsname{2.16}
\expandafter\def\csname jnum@e1.rank.m1b\endcsname{4.50}
\expandafter\def\csname jnum@e1.rank.m2grand\endcsname{1.89}
\expandafter\def\csname jnum@e1.rank.pi0\endcsname{4.33}
\expandafter\def\csname jnum@e1.wilcoxon.m2vsdroracle.a12\endcsname{0.304}
\expandafter\def\csname jnum@e1.wilcoxon.m2vsdroracle.p\endcsname{0.000528}
\expandafter\def\csname jnum@e1.wilcoxon.m2vsdrspt.a12\endcsname{0.086}
\expandafter\def\csname jnum@e1.wilcoxon.m2vsdrspt.p\endcsname{0.000000}
\expandafter\def\csname jnum@e1.wilcoxon.m2vsdrtou.a12\endcsname{0.000}
\expandafter\def\csname jnum@e1.wilcoxon.m2vsdrtou.p\endcsname{0.000000}
\expandafter\def\csname jnum@e1.wilcoxon.m2vsm1a.a12\endcsname{0.442}
\expandafter\def\csname jnum@e1.wilcoxon.m2vsm1a.p\endcsname{0.368775}
\expandafter\def\csname jnum@e1.wilcoxon.m2vsm1b.a12\endcsname{0.168}
\expandafter\def\csname jnum@e1.wilcoxon.m2vsm1b.p\endcsname{0.000000}
\expandafter\def\csname jnum@e1.wilcoxon.m2vspi0.a12\endcsname{0.165}
\expandafter\def\csname jnum@e1.wilcoxon.m2vspi0.p\endcsname{0.000000}
\expandafter\def\csname jnum@e11.charger1.eps.mean\endcsname{-0.0718}
\expandafter\def\csname jnum@e11.charger1.submod.violation\endcsname{41.1}
\expandafter\def\csname jnum@e11.charger1.utilisation\endcsname{5.6}
\expandafter\def\csname jnum@e11.charger2.eps.mean\endcsname{-0.0721}
\expandafter\def\csname jnum@e11.charger2.submod.violation\endcsname{41.2}
\expandafter\def\csname jnum@e11.charger2.utilisation\endcsname{2.8}
\expandafter\def\csname jnum@e11.charger3.eps.mean\endcsname{-0.0721}
\expandafter\def\csname jnum@e11.charger3.submod.violation\endcsname{41.2}
\expandafter\def\csname jnum@e11.charger3.utilisation\endcsname{1.9}
\expandafter\def\csname jnum@e11.lambda10.admissible\endcsname{82.5}
\expandafter\def\csname jnum@e11.lambda10.binding\endcsname{100.0}
\expandafter\def\csname jnum@e11.lambda10.margin.mean\endcsname{1.22}
\expandafter\def\csname jnum@e11.lambda125.admissible\endcsname{95.0}
\expandafter\def\csname jnum@e11.lambda125.binding\endcsname{100.0}
\expandafter\def\csname jnum@e11.lambda125.margin.mean\endcsname{1.53}
\expandafter\def\csname jnum@e11.lambda15.admissible\endcsname{100.0}
\expandafter\def\csname jnum@e11.lambda15.binding\endcsname{65.0}
\expandafter\def\csname jnum@e11.lambda15.margin.mean\endcsname{1.84}
\expandafter\def\csname jnum@e11.lambda20.admissible\endcsname{100.0}
\expandafter\def\csname jnum@e11.lambda20.binding\endcsname{15.0}
\expandafter\def\csname jnum@e11.lambda20.margin.mean\endcsname{2.45}
\expandafter\def\csname jnum@e11.penalty.eps.spread.max\endcsname{0.0000}
\expandafter\def\csname jnum@e11.penalty.gini.spread.max\endcsname{0.0000}
\expandafter\def\csname jnum@e13.norm.energy.share.mean\endcsname{9.6}
\expandafter\def\csname jnum@e13.omega.star.max\endcsname{0.0045}
\expandafter\def\csname jnum@e13.omega.star.mean\endcsname{0.0034}
\expandafter\def\csname jnum@e13.ratio.max\endcsname{488.4}
\expandafter\def\csname jnum@e13.ratio.mean\endcsname{305.6}
\expandafter\def\csname jnum@e13.ratio.min\endcsname{219.1}
\expandafter\def\csname jnum@e13.raw.energy.share.max\endcsname{0.45}
\expandafter\def\csname jnum@e13.raw.energy.share.mean\endcsname{0.34}
\expandafter\def\csname jnum@e3.calls.per.instance\endcsname{64}
\expandafter\def\csname jnum@e3.calls.total\endcsname{2560}
\expandafter\def\csname jnum@e3.core.empty.count\endcsname{39}
\expandafter\def\csname jnum@e3.core.empty.rate\endcsname{97.5}
\expandafter\def\csname jnum@e3.core.nonempty.rate\endcsname{2.5}
\expandafter\def\csname jnum@e3.eps.mean\endcsname{-0.0817}
\expandafter\def\csname jnum@e3.eps.ratio.mean\endcsname{-0.550}
\expandafter\def\csname jnum@e3.eps.ratio.median\endcsname{-0.518}
\expandafter\def\csname jnum@e3.eps.ratio.n\endcsname{38}
\expandafter\def\csname jnum@e3.eps.ratio.worst\endcsname{-2.082}
\expandafter\def\csname jnum@e3.gini.shapley.mean\endcsname{0.448}
\expandafter\def\csname jnum@e3.infeasible.coalition.rate\endcsname{0.0}
\expandafter\def\csname jnum@e3.machine.payoff.mean\endcsname{0.0093}
\expandafter\def\csname jnum@e3.mono.violation.mean\endcsname{12.61}
\expandafter\def\csname jnum@e3.ms.per.call\endcsname{92.1}
\expandafter\def\csname jnum@e3.negative.saving.count\endcsname{2}
\expandafter\def\csname jnum@e3.negative.saving.rate\endcsname{5.0}
\expandafter\def\csname jnum@e3.robot.payoff.mean\endcsname{0.0502}
\expandafter\def\csname jnum@e3.robot.share.mean\endcsname{77.4}
\expandafter\def\csname jnum@e3.robot.share.std\endcsname{30.6}
\expandafter\def\csname jnum@e3.saving.max\endcsname{66.8}
\expandafter\def\csname jnum@e3.saving.mean\endcsname{34.8}
\expandafter\def\csname jnum@e3.saving.min\endcsname{-21.6}
\expandafter\def\csname jnum@e3.saving.std\endcsname{19.2}
\expandafter\def\csname jnum@e3.seconds.total\endcsname{235.7}
\expandafter\def\csname jnum@e3.shapley.incore.rate\endcsname{0.0}
\expandafter\def\csname jnum@e3.submod.violation.max\endcsname{58.5}
\expandafter\def\csname jnum@e3.submod.violation.mean\endcsname{32.8}
\expandafter\def\csname jnum@e3.submod.violation.min\endcsname{7.7}
\expandafter\def\csname jnum@e3.v.mean\endcsname{0.1377}
\expandafter\def\csname jnum@e4.blocks.mean\endcsname{4.75}
\expandafter\def\csname jnum@e4.dhp.stable.rate\endcsname{100.0}
\expandafter\def\csname jnum@e4.maxblock.max\endcsname{3}
\expandafter\def\csname jnum@e4.maxblock.mean\endcsname{1.88}
\expandafter\def\csname jnum@e4.messages.mean\endcsname{15.5}
\expandafter\def\csname jnum@e4.partition.beats.grand.count\endcsname{2}
\expandafter\def\csname jnum@e4.partition.beats.grand.rate\endcsname{5.0}
\expandafter\def\csname jnum@e4.value.ratio.mean\endcsname{96.7}
\expandafter\def\csname jnum@e4.value.ratio.median\endcsname{100.0}
\expandafter\def\csname jnum@e4.value.ratio.min\endcsname{51.7}
\expandafter\def\csname jnum@e5.breakeven.ctime.hour\endcsname{0.31}
\expandafter\def\csname jnum@e5.breakeven.ctime.mean\endcsname{0.0052}
\expandafter\def\csname jnum@e5.front.size.mean\endcsname{3.7}
\expandafter\def\csname jnum@e5.nbs.ks.agree.rate\endcsname{97.4}
\expandafter\def\csname jnum@e5.omega.ks.mean\endcsname{0.235}
\expandafter\def\csname jnum@e5.omega.nbs.mean\endcsname{0.238}
\expandafter\def\csname jnum@e5.omega.nbs.median\endcsname{0.180}
\expandafter\def\csname jnum@e5.omega.nbs.std\endcsname{0.151}
\expandafter\def\csname jnum@exp.rules.count\endcsname{120}
\expandafter\def\csname jnum@meta.wallclock.minutes\endcsname{16.9}


\journal{Engineering Applications of Artificial Intelligence}

\begin{document}

\begin{frontmatter}

\title{Cooperative and non-cooperative game-theoretic scheduling of production machines and battery-constrained autonomous vehicles under time-of-use tariffs}

\author[1]{Kader SANOGO\corref{cor1}}\ead{ksanogo@sharjah.ac.ae}
\cortext[cor1]{Corresponding author}
\author[1]{Malek MASMOUDI}\ead{mmasmoudi@sharjah.ac.ae}
\affiliation[1]{organization={Research Institute of Science and Engineering (RISE)},
            addressline={University of Sharjah},
            city={Sharjah},
            country={United Arab Emirates}}
\author[2]{Abdelkader MEKHALEF BENHAFSSA}\ead{amekhalef@cesi.fr}
\affiliation[2]{organization={CESI LINEACT, EA 7527},
            addressline={Angouleme Campus},
            city={La Couronne},
            postcode={16400},
            country={France}}
\author[3]{M'hammed SAHNOUN}\ead{msahnoun@cesi.fr}
\affiliation[3]{organization={CESI LINEACT, EA 7527},
            addressline={Rouen Campus, S. E du Rouvray},
            city={Rouen},
            postcode={76800},
            country={France}}

\begin{abstract}
Job shops that schedule machines and battery-constrained autonomous intelligent vehicles
(AIVs) together under time-of-use (ToU) electricity tariffs face two coupled levers:
\emph{when} to draw power, and \emph{when} to recharge. Centralised optimisation of both is
computationally prohibitive at industrial scale, while existing decentralised schedulers
offer no stability guarantee and fix the makespan--energy trade-off with a hand-tuned
weight. This paper develops a decentralised mechanism with both. We formulate the problem
with the planning horizon as a hard deadline and a normalised bi-objective cost, and give
an exact constraint-programming reference that extends a validated job-shop-with-transport
solver with ToU-indexed processing cost, charger capacity and state-of-charge tracking. On
the non-cooperative side we prove that each stage game with marginal-contribution utilities
is an exact potential game whose price of stability is one, so sequential best response
converges to a pure equilibrium; the resulting mechanism reduces the global cost by
\jnum{e1.m1a.gain.mean}\,\% on average against the reference completion policy. On the
cooperative side we build a transferable-utility savings game over machines and vehicles and
show that its structure is \emph{not} the convex one the diminishing-returns intuition
suggests: a scarce shared resource that saturates makes the coalition cost supermodular, and
we exhibit a two-vehicle instance whose core is provably empty. Stability is therefore
certified per instance through the least-core radius, which we compute exactly for all
\jnum{bench.instances} benchmark instances; the core is empty on
\jnum{e3.core.empty.rate}\,\% of them, with a mean radius of
\jnum{e3.eps.ratio.mean} of total savings. Decentralised merge--split coalition formation
recovers \jnum{e4.value.ratio.mean}\,\% of the grand-coalition savings using coalitions of
at most \jnum{e4.maxblock.max} agents (mean \jnum{e4.maxblock.mean}) and
\jnum{e4.messages.mean} bilateral messages, and
Nash bargaining between the production and logistics blocs selects a trade-off weight of
\jnum{e5.omega.nbs.mean} rather than leaving it to be assumed. Code, benchmark and every
reported number are released.
\end{abstract}

\begin{graphicalabstract}
\end{graphicalabstract}

\begin{highlights}
\item Decentralised game-theoretic scheduling of machines and battery-constrained AIVs.
\item Each stage game is an exact potential game whose price of stability equals one.
\item Resource congestion makes coalition costs supermodular and can empty the core.
\item Least-core radius certifies stability where no core-existence theorem can.
\item Bargaining endogenises the makespan-energy weight; code and benchmark are released.
\end{highlights}

\begin{keyword}
Job shop scheduling with transportation \sep Cooperative game theory \sep Least core \sep
Time-of-use tariffs \sep Autonomous intelligent vehicles \sep Decentralised optimisation
\end{keyword}

\end{frontmatter}

\section{Introduction}
\label{sec:intro}

Modern job shops are asked to be fast, flexible and cheap to run at the same time. Under
Industry~5.0 the third of those has acquired a specific meaning: electricity is no longer
billed at a flat rate but under time-of-use (ToU) tariffs that price the same kilowatt-hour
differently depending on when it is drawn~\citep{7349258, su16062443}. At the same time,
material handling inside the shop has moved from fixed conveyors to fleets of autonomous
intelligent vehicles (AIVs) whose finite batteries must themselves be recharged --- an
activity that is both a demand on the tariff and a period of unavailability for
transport~\citep{DESTOUET2023155, en13184948}.

These two developments interact in a way that neither classical job-shop scheduling nor its
transport-aware extension captures. Deferring an operation into an off-peak band saves money
but consumes deadline slack; recharging a vehicle early is cheap but removes it from service
when the shop is busiest; and because the charging station is a shared resource, one
vehicle's cheap slot is another's queue. The decision problem is a Job Shop Scheduling
Problem with Transportation (JSSPT) extended along two coupled axes --- an economic axis
(tariff-dependent cost) and a resource-feasibility axis (state of charge) --- with a hard
completion deadline.

Centralised optimisation of the combined problem does not scale. Section~\ref{sec:results}
reports that an exact constraint-programming model closes the pure-makespan variant of a
four-machine benchmark instance in under two seconds but cannot close the ToU-and-battery
variant of the same instance in \jnum{e2.budget.seconds} seconds, because pricing energy
against the tariff destroys the makespan-optimal plateau on which the classical model
relies. That is the concrete form of the scalability argument usually made in the abstract,
and it is what motivates decentralisation here.

Decentralised scheduling, however, buys tractability with a question: if machines and
vehicles each decide for themselves, what stops the resulting schedule from being both
inefficient and unstable? Potential games answer half of that question --- they guarantee
that self-interested best response converges --- but they say nothing about whether the
agents would \emph{stay} in the arrangement if they could renegotiate, nor about how the
savings should be divided. Those are cooperative questions, and they are the ones this paper
takes up.

\paragraph{Delta with respect to our own prior work.}
We have previously applied potential games with marginal-contribution utilities to job-shop
scheduling with transport in shared human--robot environments~\citep{SANOGO2025111366}, and
to transport-integrated scheduling more broadly~\citep{SANOGO2025103060, SANOGO2023209}. The
non-cooperative mechanism of Section~\ref{sec:m1} builds directly on that line and we do not
claim it as new. What is new here is: (i) time-of-use tariffs and battery/charging decisions
as first-class decision variables, which change what the agents are optimising and introduce
the shared charging resource; (ii) an entire \emph{cooperative} layer --- a
transferable-utility savings game over machines and vehicles, its stability certification,
and a bargaining-derived trade-off weight --- which has no counterpart in the prior work;
(iii) a corrected equilibrium theory, in which the potential property is proved at the scope
at which it holds and the optimality claim is replaced by a price-of-stability result and a
measured price of anarchy; and (iv) an exact constraint-programming reference for the
combined problem. No text is reused from the earlier papers.

\paragraph{Contributions.}
\begin{itemize}
  \item \textbf{A cooperative scheduling game with certified stability.} We define a
  characteristic cost function for an integrated production--transport system under ToU
  tariffs, and certify stability \emph{per instance} through the least-core radius
  (Theorem~\ref{thm:leastcore}) rather than through a core-existence theorem that does not
  hold here.
  \item \textbf{A structural negative result.} We show by explicit construction
  (Proposition~\ref{prop:n5}) that congestion in a scarce shared resource makes the
  coalition cost supermodular, so the savings game is concave and the core can be empty ---
  and we verify that the obstruction is present throughout the benchmark, not only in the
  constructed instance.
  \item \textbf{An endogenised trade-off weight.} The makespan--energy weight is derived
  from a Nash (and Kalai--Smorodinsky) bargaining solution between the production and
  logistics blocs, with the non-cooperative equilibrium front supplying the disagreement
  point (Section~\ref{sec:layer3}).
  \item \textbf{Decentralised coalition-structure generation} by merge and split over a
  hedonic preference induced by within-coalition Shapley payoffs, converging to a
  $\mathbb{D}_{hp}$-stable partition with a reported communication cost.
  \item \textbf{Corrected equilibrium theory}: an exact stage potential
  (Theorem~\ref{thm:t1}), an explicit statement of the scope at which it fails
  (Section~\ref{sec:t2}), price of stability equal to one (Theorem~\ref{thm:t3}), and an
  empirical price of anarchy measured against an exact bound.
  \item \textbf{An exact CP-SAT reference} for the combined problem, built as an extension
  of a validated JSSPT solver and regression-tested against it.
  \item \textbf{A released artefact}: simulator, benchmark files with regenerated tariff
  profiles, all solution concepts, and a registry that maps every number in this paper to
  the run that produced it.
\end{itemize}

Section~\ref{sec:rwork} positions the work. Section~\ref{sec:problem} states the problem,
including the units contract and the deadline. Section~\ref{sec:exact} gives the exact
reference. Section~\ref{sec:m1} develops the non-cooperative game and
Section~\ref{sec:m2} the cooperative one. Sections~\ref{sec:exp}--\ref{sec:results} report
the experiments, Section~\ref{sec:insights} the managerial reading, and
Section~\ref{sec:conc} concludes.

\section{Related work}
\label{sec:rwork}

\subsection{Energy-aware scheduling under time-of-use tariffs}
Energy-aware production scheduling divides into approaches that reduce total consumption ---
by speed scaling, machine on/off control or turning off auxiliary equipment --- and
approaches that reduce \emph{cost} by displacing a fixed consumption in
time~\citep{7349258, MANSOURI2016772, ZHANG2017665}. Under a ToU tariff with constant
processing power and idle power excluded, total machine energy is schedule-invariant, so the
entire lever is temporal displacement. That is the framing adopted here, and it is the one
for which ToU pricing is the right instrument~\citep{su14106264, app12031491}. Recent work
extends this to distributed and heterogeneous shops under non-identical
tariffs~\citep{WANG2024777, LENG2025528, ZHAO2025420}, but almost always without a transport
resource: material handling is assumed instantaneous or absorbed into processing times.

\subsection{Job-shop scheduling with transport and battery-constrained vehicles}
The JSSPT has been studied since~\citet{bilge1995}, whose ten job sets and four layouts
remain the reference benchmark and are used here. Battery constraints entered the literature
much later, largely through AGV/AIV fleet management~\citep{DESTOUET2023155, DESTOUET2024110419,
QIN2023242}. Two modelling choices recur and both matter for what follows. First, battery
feasibility is often imposed as an \emph{aggregate} energy balance over a vehicle's assigned
tasks, which ignores ordering: a vehicle can satisfy the total-energy inequality and still
cross its state-of-charge floor midway through its sequence, so such a model is a relaxation
and its objective a bound. Second, the charging station's own capacity is frequently left
unconstrained, which removes precisely the congestion that Section~\ref{sec:m2} shows to be
structurally decisive. Our exact model in Section~\ref{sec:exact} imposes both.

\subsection{Game-theoretic and multi-agent scheduling}
Decentralised scheduling as a game has a long tradition, with potential
games~\citep{monderer1996potential} and distributed welfare games~\citep{marden2009cooperative}
providing the standard machinery: assign each agent a marginal-contribution utility
(\citealp{9241495, chapman2009}) and the induced game admits a pure Nash equilibrium
reachable by best response. Equilibrium quality is then discussed through the price of
anarchy~\citep{koutsoupias2009worst} and the price of stability~\citep{anshelevich2008price}.
Two cautions apply and are respected here. The potential property holds for a stage game
whose payoff is a function of the current joint action, not for a payoff that depends on an
entire downstream trajectory. And bounds derived for \emph{maximisation} of a monotone
submodular welfare~\citep{1181966} do not transfer to a cost objective by changing sign:
negating a monotone submodular function gives a non-increasing supermodular one, violating
both hypotheses.

\subsection{Cooperative game theory in operations and manufacturing}
Cooperative solution concepts --- the core, the Shapley value~\citep{shapley1971cores}, the
nucleolus~\citep{schmeidler1969nucleolus} --- are standard in collaborative logistics and
cost allocation, and the classical result is that a \emph{convex} game has a non-empty core
containing the Shapley value~\citep{shapley1971cores}; the converse direction is due
to~\citet{ichiishi1981super} and the two are frequently conflated. When the core is empty,
the least core~\citep{maschler1979geometric} still quantifies how far from stable the best
allocation is, and it is always computable. Coalition-structure generation by merge and
split~\citep{apt2009stable, saad2009coalitional} keeps coalitions small enough to be
negotiated locally. To our knowledge no prior work analyses core stability for a joint
machine/AIV energy-aware scheduling game; that gap is what this paper addresses, and the
answer turns out to be a negative structural result plus a certification procedure rather
than an existence theorem.

\subsection{Positioning}
Table~\ref{tab:positioning} places the closest work along the seven features that matter
here. No prior study combines transport as a scheduled resource, battery state of charge,
charging as a decision, a charger capacity constraint, a hard deadline, a decentralised
mechanism, and a cooperative stability analysis. The first six have each appeared; the
seventh has not appeared at all in this problem class, and it is the one that changes the
conclusions.

\begin{table}[t]
\centering
\small
\caption{Positioning against the closest literature. \checkmark: present; --: absent.}
\label{tab:positioning}
\begin{tabular}{@{}lccccccc@{}}
\toprule
Work & Transport & Battery & Charging & Charger & Deadline & Decentral. & Cooperative \\
     & resource   & SoC     & decision & capacity &          & mechanism  & stability \\
\midrule
\citet{bilge1995}            & \checkmark & -- & -- & -- & -- & -- & -- \\
\citet{7349258}              & -- & -- & -- & -- & \checkmark & -- & -- \\
\citet{ZHANG2017665}         & -- & -- & -- & -- & -- & -- & -- \\
\citet{MANSOURI2016772}      & -- & -- & -- & -- & \checkmark & -- & -- \\
\citet{DESTOUET2023155}      & \checkmark & \checkmark & \checkmark & -- & -- & -- & -- \\
\citet{QIN2023242}           & \checkmark & \checkmark & \checkmark & -- & -- & \checkmark & -- \\
\citet{WANG2024777}          & -- & -- & -- & -- & -- & \checkmark & -- \\
\citet{LENG2025528}          & -- & -- & -- & -- & -- & \checkmark & -- \\
\citet{SANOGO2025111366}     & \checkmark & -- & -- & -- & -- & \checkmark & -- \\
\citet{SANOGO2025103060}     & \checkmark & -- & -- & -- & -- & \checkmark & -- \\
\textbf{This paper}          & \checkmark & \checkmark & \checkmark & \checkmark & \checkmark & \checkmark & \checkmark \\
\bottomrule
\end{tabular}
\end{table}

\section{Problem statement}
\label{sec:problem}

\subsection{Sets, tasks and decisions}
A set of jobs $\mathcal{J}=\{J_1,\dots,J_n\}$ is processed on machines
$\mathcal{M}=\{M_1,\dots,M_{|\mathcal{M}|}\}$ and transported by a fleet of single-load AIVs
$\mathcal{V}=\{V_1,\dots,V_{|\mathcal{V}|}\}$. Job $J_i$ is an ordered chain of operations
$O_{i,1},\dots,O_{i,m_i}$; operation $O_{i,k}$ requires machine $\mu(i,k)$ for $p_{i,k}$
time units without pre-emption. Locations $\mathcal{L}$ comprise the machines, the
load/unload station $LU$ and the charging station $CH$, with travel times $t_{a,b}$.

Every operation is preceded by a transport leg, and every job ends with a return leg to
$LU$, so the manufacturing sequence of $J_i$ is
$(T_{i,1}, O_{i,1}, T_{i,2}, \dots, O_{i,m_i}, T_{i,m_i+1})$ and the makespan $C_{\max}$ is
measured to the \emph{last return to $LU$}. Transport task $T_{i,k}$ moves the part from
$\mathrm{prev}(i,k)$ to $\mathrm{dest}(i,k)$ in $\tau_{i,k}=t_{\mathrm{prev}(i,k),\mathrm{dest}(i,k)}$.

Decisions are event-driven and taken by the resources themselves. An idle machine chooses
which buffered job to start and, from a finite ToU-aligned grid $\mathcal{D}$, how long to
defer it. A free vehicle chooses a precedence-ready transport task or a charging visit,
again with a deferral from $\mathcal{D}$. Decisions occur only at events --- a vehicle
becoming free, a machine becoming free, a part arriving, a charger being released --- so the
number of decisions scales with system activity rather than with horizon length.

\subsection{Tariff and horizon}
The horizon is partitioned into intervals $\mathcal{H}$ with prices $c_h$. Following the
convention established for tariff-aligned scheduling, the horizon length is
\begin{equation}
H \;=\; \lambda\left(\frac{1}{|\mathcal{M}|}\sum_{i,k} p_{i,k}
        \;+\; \frac{1}{|\mathcal{V}|}\sum_{i,k}\tau_{i,k}\right),
\label{eq:horizon}
\end{equation}
with $\lambda=2$ in the published setting. Six periods tile $[0,H)$ at prices
$(0.14, 0.22, 0.14, 0.22, 0.14, 0.06)$~\euro/kWh in proportions $(3,5,3,3,2,8)/24$ of $H$.
Interval lengths are obtained by rounding $H f_k$ half away from zero for the first five
periods, the sixth absorbing the remainder, so that the profile tiles $[0,H)$ exactly and
the last period never ends after $H$.

\subsection{The horizon is a hard deadline}
\label{sec:deadline}
$H$ is the maximum permitted completion time: no schedule may finish after it.
\begin{equation}
C_{\max} \;\le\; H .
\label{eq:deadline}
\end{equation}
Stating this explicitly has four consequences that are design obligations rather than
algebra, and each is discharged in this paper.

\emph{(i) Instances can be infeasible.} $H$ is generated by~\eqref{eq:horizon} and is not
guaranteed to admit any feasible schedule, so each instance carries a certified
\emph{feasibility margin} $H/C_{\max}^{LB}$ against a valid lower bound. All
\jnum{bench.instances} benchmark instances are admissible at $\lambda=2$
(Section~\ref{sec:res-lambda}), but the margin is loose: $\lambda$ is therefore a
feasibility parameter, not only a tariff-shaping one, and its ablation is load-bearing.

\emph{(ii) Decentralised dynamics can deadlock.} A myopically cheap decision --- deferring a
start into an off-peak band, or sending a vehicle to charge --- can leave no feasible
completion. Every method here checks a lower bound on remaining work against the time left
before committing, and falls back to a deadline-directed policy when the bound is violated.
The \emph{deadlock rate}, the fraction of episodes needing the fallback, is reported as a
first-class metric: a method that meets the deadline only by falling back is not a
scheduler.

\emph{(iii) The deadline supplies the upper normalisation anchor.} $H$ is exactly the value
of $C_{\max}$ a feasible schedule may not exceed, which is what makes the normalisation of
Section~\ref{sec:objective} principled rather than arbitrary.

\emph{(iv) The coalition cost could become infinite.} A small coalition may have no
deadline-feasible completion, which would make the characteristic function non-real-valued
and remove the Shapley value, the core programme and the least-core radius on exactly the
instances where coordination matters most. Section~\ref{sec:m2} adopts a finite penalised
form.

\subsection{Units contract}
\label{sec:units}
Battery capacities are quoted in ampere-hours while tariffs are quoted per kilowatt-hour, so
a nominal bus voltage must be declared before either can be converted into the other. Three
reconciliations follow and are given in Table~\ref{tab:units}.

First, a $100$~Ah pack at $U_{\mathrm{nom}}=48$~V stores $4.8$~kWh. We assume fast charging
takes the pack from $0\,\%$ to $80\,\%$ of capacity in $25$~minutes with the state of charge
growing linearly, which fixes the charging rate at $3.2$~Ah/min, hence $16$~Ah per
five-minute block and an implied charger power of $9.216$~kW.

Second, the depletion rates $1/3/7$ are read as \textbf{mAh per second}, i.e.\ $60/180/420$
mAh/min while idle, travelling empty and travelling loaded. This matters: read literally as
milliamperes per second, $7$ would drain $0.043$~Ah over the longest horizon in the
benchmark, the state-of-charge floor could never bind, and the battery dimension of the
problem would be inert. At $7$~mAh/s the usable $60$~Ah window drains in about $143$
minutes, and the floor binds on \jnum{units.floor.binding.rate}\,\% of benchmark instances
under a purely reactive charging rule.

Third, the battery-feasibility predicate is expressed entirely in mAh, so the floor enters
as $20\,000$~mAh:
\begin{equation}
B_{\mathrm{needed}}(a) = B^{\min}
 + q^{e}\,t_{p,i} + q^{\mathrm{idle}}\,[\,t^{r}_{i,k-1}-t_{p,i}\,]^{+}
 + q^{\ell}\,t_{i,d} + q^{e}\,t_{d,CH},
\label{eq:bneeded}
\end{equation}
where $q^{e}, q^{\mathrm{idle}}, q^{\ell}$ are the empty, idle and loaded rates. The final
term is what makes the predicate safe: after completing the task the vehicle must still be
able to reach the charger without crossing the floor. The $[20\,\%,80\,\%]$ window is
enforced at both ends; charging stops at the ceiling.

Vehicle traction energy is charged for at \emph{charging} time, when the grid draw actually
occurs, so a transport action has zero direct cost while fully depleting the battery. This is
internally consistent but easy to misread, and is stated here explicitly.

\begin{table}[t]
\centering
\small
\caption{Units contract. Derived quantities are marked; the charger power is derived from
the recharge block rather than assumed, because the two cannot both be free parameters.}
\label{tab:units}
\begin{tabular}{@{}llll@{}}
\toprule
Quantity & Symbol & Value & Note \\
\midrule
Time                       & --                 & minutes                & all data integral \\
Energy                     & --                 & kWh                    & \\
Money                      & --                 & \euro                  & tariff in \euro/kWh \\
Charge                     & --                 & mAh                    & \\
Nominal bus voltage        & $U_{\mathrm{nom}}$ & 48 V                   & declared \\
Battery capacity           & $B^{\max}$         & 100 Ah $=100\,000$ mAh & $=4.8$ kWh \\
Operating window           & $[B^{\min},B^{\mathrm{cap}}]$ & $[20\,000, 80\,000]$ mAh & both ends enforced \\
Depletion, idle/empty/loaded & $q^{\mathrm{idle}},q^{e},q^{\ell}$ & $60/180/420$ mAh/min & read as mAh/s \\
Fast charge                & --                 & $0\to80\,\%$ in 25 min & assumed, linear \\
Charging rate              & --                 & 3.2 Ah/min             & derived \\
Charging block             & --                 & 5 min $\to$ 16 Ah      & derived \\
Charger power              & $P^{\mathrm{ch}}$  & 9.216 kW               & derived \\
Machine power              & $P^{\mathrm{mach}}$& 2 kW while processing  & \citet{VISWANATHAN20221409} \\
Charger capacity           & $K_{CH}$           & 1                      & unary by default \\
\bottomrule
\end{tabular}
\end{table}

\subsection{Objective}
\label{sec:objective}
Both energy streams are priced against the tariff:
\begin{equation}
E_{\mathrm{cost}} \;=\;
\sum_{(i,k)}\int_{S_{i,k}}^{S_{i,k}+p_{i,k}} P^{\mathrm{mach}}\,c(t)\,\mathrm{d}t
\;+\;
\sum_{r\in\mathcal{V}}\ \sum_{\text{charging blocks of } r}\int P^{\mathrm{ch}}\,c(t)\,\mathrm{d}t .
\label{eq:ecost}
\end{equation}
Both terms are piecewise-constant integrals over the tariff partition and are computed
exactly. Including machine processing energy is not optional: a model that prices only
charging optimises a different quantity from one that prices both, and no comparison between
them is meaningful.

\paragraph{Commensurability.} A weighted sum $\omega C_{\max} + (1-\omega)E_{\mathrm{cost}}$
adds \emph{minutes to euros}. On this benchmark the two raw terms differ by a factor of
\jnum{e13.ratio.min}--\jnum{e13.ratio.max}, so at $\omega=0.5$ the energy term contributes
a mean of \jnum{e13.raw.energy.share.mean}\,\% of the objective and the two balance only near
$\omega\approx\jnum{e13.omega.star.mean}$. Any sweep over $\omega\in[0,1]$ on the raw
objective is therefore a makespan study with an energy label. We normalise first, using
anchors that are closed-form constants of the instance and require no solver call:
\begin{equation}
\widehat{C}_{\max}=\frac{C_{\max}-C_{\max}^{LB}}{H-C_{\max}^{LB}},\qquad
\widehat{E}=\frac{E_{\mathrm{cost}}-E^{LB}}{E^{UB}-E^{LB}},\qquad
\Phi = \omega\widehat{C}_{\max} + (1-\omega)\widehat{E}.
\label{eq:phi}
\end{equation}
$C_{\max}^{LB}$ is the maximum of the job critical path including transport, a one-machine
head/tail bound, and a fleet work bound; $E^{LB}$ and $E^{UB}$ price the same reference
kilowatt-hours at the cheapest and dearest tariff, the reference being the schedule-invariant
processing energy plus a charging quantity bracketed between the fleet's unavoidable
discharge and the fastest discharge the horizon permits. After normalisation the energy term
contributes a mean of \jnum{e13.norm.energy.share.mean}\,\% at $\omega=0.5$, and $\omega$
does what the notation says it does.

\begin{assumption}[Anchor invariance]
\label{ass:anchors}
The anchors are computed once per instance, before any episode begins, and never updated.
Recomputing them from running best-known values would make $\Phi$ a function of history
rather than of the joint action profile, and the exact-potential property of
Theorem~\ref{thm:t1} would fail with it. This is asserted in the released test suite.
\end{assumption}

Deadline violations are charged a finite penalty $M\max(0,\widehat{C}_{\max}-1)$, zero on
feasible outcomes. Finiteness is what keeps every cooperative solution concept defined
(Section~\ref{sec:m2}); $M=10$ by default and the sensitivity to it is ablated.

\section{Exact reference model}
\label{sec:exact}

The reference model \textsc{ex-cp} is a constraint-programming model built with
CP-SAT~\citep{ortools}. It is an \emph{extension} of a validated job-shop-with-transport
solver rather than a reimplementation: the skeleton --- interval variables per operation,
\texttt{NoOverlap} per machine, a transport leg before every operation, a final return leg,
exactly one vehicle per transport, and sequence-dependent empty travel between consecutive
transports --- is retained unchanged, and five elements are added.

\begin{enumerate}
\item \textbf{ToU-indexed processing cost.} For each operation and each tariff interval $h$
an overlap variable $o_{i,k,h}=\big[\min(C_{i,k},T_h)-\max(S_{i,k},T_{h-1})\big]^{+}$ is
built from min/max equalities, and the cost contribution is
$\sum_h c_h P^{\mathrm{mach}} o_{i,k,h}/60$. The same construction prices charging.
\item \textbf{Charging as optional intervals}, one per potential visit per vehicle, with a
symmetry break ordering the slots.
\item \textbf{Charger capacity} through \texttt{Cumulative} over all charging intervals with
capacity $K_{CH}$.
\item \textbf{Exact prefix state of charge.} This is the one constraint block that had to be
\emph{ported} rather than reused. The original sequences transports on a vehicle with
pairwise big-$M$ disjunctions, which fix relative order but never adjacency, and therefore
cannot express a prefix state such as the state of charge --- nor interleave charging visits
with transports. \textsc{ex-cp} replaces that block with one \texttt{AddCircuit} per vehicle
over $\{\text{depot}\}\cup\text{transports}\cup\text{that vehicle's charging slots}$. Circuit
arcs give adjacency directly, so the empty leg, the idle wait and the recursion
\begin{equation}
\mathrm{soc}_b = \min\!\big(B^{\mathrm{cap}},\;
\mathrm{soc}_a - q^{e}t_{\mathrm{drop}(a),\mathrm{pick}(b)} - q^{\mathrm{idle}}w_{ab}
- q^{\ell}\tau_b + g_b\big)
\label{eq:socrec}
\end{equation}
are all exact along the route, with $\mathrm{soc}_b \ge B^{\min}$ enforced at every node.
An aggregate energy balance, by contrast, is a relaxation: a vehicle can satisfy the
total-energy inequality and still cross the floor midway through its sequence.
\item \textbf{The deadline} $C_{\max}\le H$ as a hard constraint, which also tightens the
domain of every end variable.
\end{enumerate}

The objective is the same normalised $\Phi$ of~\eqref{eq:phi}, which is affine in the two
integer quantities $C_{\max}$ and $E_{\mathrm{cost}}$ and therefore exactly representable
once the coefficients are scaled to integers. Times are integer minutes, charge is integer
mAh and money is integer micro-euros, so the solver never leaves integer arithmetic.

\paragraph{Regression against the original.} Setting every tariff interval to the same price
and removing both the deadline and the battery makes the energy term schedule-invariant, so
\textsc{ex-cp} must return \emph{exactly} the makespan the original solver returns. It does,
on every instance tested, and this is a permanent test in the released artefact rather than a
one-off check. It is the cheapest correctness guarantee available here: any discrepancy
would be a bug in the adaptation, located before any tariff result exists.

\section{$M1$: the non-cooperative stage game}
\label{sec:m1}

\subsection{Utilities, in one convention}
Players are $\mathcal{P}=\mathcal{M}\cup\mathcal{V}$. Fix a state $s_q$ and let
$\widehat{\Phi}(s_q,\alpha)$ be the value of $\Phi$ obtained by committing the joint action
$\alpha$ through a fixed deterministic execution rule and completing the episode with a fixed
deterministic reference policy $\pi_0$. Define
\begin{equation}
W(\alpha) \;=\; \widehat{\Phi}(s_q,\alpha^{0}) - \widehat{\Phi}(s_q,\alpha),
\qquad
u_p(\alpha) \;=\; W(\alpha) - W(a_p^{0},\alpha_{-p})
\;=\; \widehat{\Phi}(a_p^{0},\alpha_{-p}) - \widehat{\Phi}(\alpha),
\label{eq:utility}
\end{equation}
where $\alpha^{0}$ is the all-null profile. Welfare $W$ is a \emph{cost reduction}, agent
utility is a marginal contribution to it, and \textbf{every agent maximises}. Stating the
convention once and using it everywhere matters: defining $u_p$ as a cost increase in one
place and a cost reduction in another, and then minimising one and maximising the other,
makes every subsequent equation ambiguous.

\paragraph{The null action is a baseline, not a strategy.}
$a_p^{0}$ is the reference against which marginal contribution is measured, but it is removed
from the \emph{choice} set whenever a player has productive work available. Two reasons, both
substantive. It is redundant: deferral is already a first-class action through the grid
$\mathcal{D}$, so a machine that wants to wait for the off-peak band says so with a delay
rather than by refusing to act. And leaving it in is degenerate: $u_p(a_p^{0})=0$
identically, while any real action must beat a completion policy that will do the same work a
moment later, so idling weakly dominates and best-response dynamics stall the shop at the
first event at which every clock has caught up. A vehicle whose only options are charging
keeps the null action, because charging genuinely is optional.

\subsection{Exact stage potential}

\begin{theorem}[Exact stage potential]
\label{thm:t1}
For any state $s_q$, any player $p$, any $a_p,a_p'\in\mathcal{A}_p(s_q)$ and any
$\alpha_{-p}$,
\[
u_p(a_p,\alpha_{-p}) - u_p(a_p',\alpha_{-p}) \;=\; W(a_p,\alpha_{-p}) - W(a_p',\alpha_{-p}).
\]
The stage game is therefore an exact potential game with potential $W$. Since the action
sets are finite, the finite improvement property holds, and \emph{sequential} best-response
dynamics reach a pure Nash equilibrium of the stage game in finitely many
steps~\citep{monderer1996potential}.
\end{theorem}

\begin{proof}
By~\eqref{eq:utility}, $u_p(a_p,\alpha_{-p}) = \widehat{\Phi}(a_p^{0},\alpha_{-p}) -
\widehat{\Phi}(a_p,\alpha_{-p})$ and likewise for $a_p'$. Subtracting, the term
$\widehat{\Phi}(a_p^{0},\alpha_{-p})$ does not depend on $p$'s own action and cancels,
leaving $\widehat{\Phi}(a_p',\alpha_{-p})-\widehat{\Phi}(a_p,\alpha_{-p})$, which is exactly
$W(a_p,\alpha_{-p})-W(a_p',\alpha_{-p})$. Finiteness of $\mathcal{A}_p$ follows from the
finite delay grid $\mathcal{D}$ and the finite task set. The finite improvement property then
gives termination provided deviations are sequential; genuinely synchronous best response can
cycle even in a potential game.
\end{proof}

Two hypotheses in Theorem~\ref{thm:t1} carry real content and are easy to lose.
\emph{Finiteness} of the action sets is why the machine action space must be
$(J_i, d)$ with $d\in\mathcal{D}$ finite, and not a continuous start time. And
$\widehat{\Phi}$ must be a function of $(s_q,\alpha)$ \emph{alone}, which is why the
completion policy $\pi_0$ and the conflict-resolution rule are both fixed and deterministic:
they are what turn a trajectory-dependent quantity into a well-defined stage payoff.

\subsection{What the potential property does not give}
\label{sec:t2}
The multi-stage game is a \emph{sequence} of stage potential games. Convergence at every
stage does not imply that the resulting trajectory is optimal, and the exact-potential
identity does not survive if $\widehat{\Phi}$ is allowed to depend on how downstream players
will re-optimise. We therefore make no trajectory-level optimality claim and quantify the
loss empirically instead (Section~\ref{sec:results}).

\begin{theorem}[Price of stability]
\label{thm:t3}
For each stage game, a global minimiser of $\widehat{\Phi}$ over the joint action space is a
pure Nash equilibrium. Hence $\mathrm{PoS}=1$: an optimal equilibrium exists at every stage.
\end{theorem}

\begin{proof}
Let $\alpha^\star \in \arg\min_\alpha \widehat{\Phi}(s_q,\alpha)$ and suppose some player $p$
had a profitable deviation $a_p$, i.e.\ $u_p(a_p,\alpha^\star_{-p}) >
u_p(\alpha^\star_p,\alpha^\star_{-p})$. By~\eqref{eq:utility} this is equivalent to
$\widehat{\Phi}(a_p,\alpha^\star_{-p}) < \widehat{\Phi}(\alpha^\star)$, contradicting
minimality. Hence no profitable unilateral deviation exists.
\end{proof}

Theorem~\ref{thm:t3} is the correct version of the informal claim that ``the optimal solution
is always an equilibrium''. It should not be confused with an upper bound on the cost of an
\emph{arbitrary} equilibrium. In particular, a bound of the form
$\tfrac12\Phi(\alpha^{\mathrm{opt}}) \le \Phi(\alpha^\star) \le \Phi(\alpha^{\mathrm{opt}})$
cannot hold for a minimisation: its right-hand inequality asserts that every equilibrium
costs no more than the optimum, which would make the optimisation vacuous. The valid-utility
theorem it derives from is a \emph{maximisation} result requiring a normalised, non-decreasing
and submodular welfare~\citep{1181966}, and the map from welfare to cost is not a sign flip:
negating a monotone submodular function yields a non-increasing supermodular one, violating
both hypotheses. The quantity that \emph{is} meaningful for an arbitrary equilibrium is the
price of anarchy, $\mathrm{PoA}=\max_{\alpha^\star\in NE}\Phi(\alpha^\star)/\Phi(\alpha^{\mathrm{opt}})\ge 1$,
which must be measured; Section~\ref{sec:res-exact} reports why it cannot yet be measured
here and what is reported instead.

\subsection{Algorithms}
\textbf{M1a (best response).} At each event, the active players are ordered by a random
permutation and updated round-robin: each in turn adopts
$\arg\max_{a\in\mathcal{A}_p} u_p(a,\alpha_{-p})$ until no player improves. Updates are
sequential by construction. Instrumentation records iterations to convergence, the number of
improvement steps, whether the iteration cap binds, and an $\varepsilon$-Nash certificate ---
the largest unilateral improvement still available at the committed profile.

\textbf{M1b (log-linear learning).} Replacing the arg-max with a Gibbs choice
$\Pr(a_p)\propto\exp(u_p(a_p,\alpha_{-p})/\theta)$ and annealing $\theta$ makes the induced
chain concentrate, as $\theta\to0$, on potential maximisers~\citep{blume1993statistical}. It
is thus an \emph{equilibrium-selection} device targeting the best equilibrium rather than an
arbitrary one. A final greedy sweep anneals to zero exactly. M1b is a selection mechanism,
not an equilibrium solver: because its final sweep is a single sequential pass, the committed
profile is an $\varepsilon$-equilibrium with small but generally non-zero $\varepsilon$.

Each utility evaluation costs one deterministic rollout, and the null baseline
$\widehat{\Phi}(a_p^{0},\alpha_{-p})$ is shared across a player's whole candidate set, which
halves the rollout count of every sweep.

\section{$M2$: the cooperative game}
\label{sec:m2}

The cooperative model has three interacting layers. Layer~2 carries the stability analysis
and is presented first, because a negative structural result found there determines how the
other two must be framed. Layer~1 makes the approach decentralised; Layer~3 removes the
arbitrary weight $\omega$.

\subsection{Layer 2: the savings game}
\label{sec:layer2}

\begin{definition}[Characteristic cost]
\label{def:c}
For $S\subseteq N=\mathcal{M}\cup\mathcal{V}$, let $c(S)$ be the value of $\Phi$ realised
when the members of $S$ coordinate --- each running best response on its
marginal-contribution utility --- while every agent in $N\setminus S$ follows the fixed
reference policy $\pi_0$, all evaluated on the shared simulator. In particular
$c(\emptyset)=\Phi(\pi_0 \text{ everywhere})$.
\end{definition}

Two properties of this ``coordinate within $S$, default outside'' construction must be stated
rather than left implicit. It \emph{fixes the externalities} from the complement, which is
what licenses a partition-function-free transferable-utility game here. And $c(S)$ is a
\emph{whole-system} cost for every $S$, singletons included.

\begin{definition}[Savings game]
\label{def:v}
$v(S) = c(\emptyset)-c(S)$, with $v(\emptyset)=0$.
\end{definition}

The alternative $v(S)=\sum_{i\in S}c(\{i\}) - c(S)$ must \emph{not} be used here. It
presupposes that $c(\{i\})$ is a standalone cost, whereas under Definition~\ref{def:c} it is
a whole-system cost; the sum then double-counts roughly $(|S|-1)c(\emptyset)$, gives
$v(\emptyset)=-c(\emptyset)\neq0$, forces $v(\{i\})=0$ for every $i$ by construction, and
leaves $v(N)$ measuring a quantity that is not the total savings --- which the Shapley
efficiency condition would then allocate.

\paragraph{Keeping $c$ real-valued.} Under the deadline~\eqref{eq:deadline} a small coalition
coordinating against a $\pi_0$ complement may have no feasible completion, which would give
$c(S)=+\infty$. A characteristic function that is not real-valued has no Shapley value, no
core programme and no least-core radius. The penalty term of Section~\ref{sec:objective}
resolves this: $c$ stays finite on all of $2^{N}$, the penalty is zero wherever the deadline
is met, and the fraction of coalitions that miss it is reported per instance
(Section~\ref{sec:res-coop}).

\paragraph{An economy in the construction.} Because the stage game is an exact potential game
with potential $W$ (Theorem~\ref{thm:t1}), best response on the \emph{individual} marginal
contribution is coordinate descent on the coalition's contribution to the global cost. A
coalition therefore needs no separate objective from its members' --- a structural
convenience, not a modelling shortcut. Following the cost discipline of this problem, $c(S)$
is evaluated by a \emph{single sequential best-response sweep} rather than a full
optimisation; results are memoised across coalitions, and both the call count and the
wall-clock are reported so that a truncated budget cannot read as complete certification.

\subsection{The congestion obstruction}
\label{sec:n5}

The intuition that coordination exhibits diminishing returns --- and that $c$ is therefore
submodular, $v$ convex, and the core non-empty with the Shapley value inside it
\citep{shapley1971cores} --- is the natural hypothesis. It is false here.

\begin{proposition}[Congestion breaks submodularity]
\label{prop:n5}
$c$ is not submodular in general. Consequently the savings game is not convex in general,
and the core can be empty.
\end{proposition}

\begin{proof}[Proof by construction]
Take two vehicles $V_1,V_2$, two single-operation jobs, a charging station of unary capacity,
and a tariff with exactly one cheap window whose length equals one charging block. Set the
initial state of charge so that each vehicle must charge exactly once, and set
$\omega=0$ with machine power zero, so that all cost is charging cost and the only decision
is \emph{when} each vehicle charges. Then:
\[
c(\emptyset)=\jnum{n5.c.empty},\quad
c(\{V_1\})=c(\{V_2\})=\jnum{n5.c.single},\quad
c(\{V_1,V_2\})=\jnum{n5.c.pair},
\]
so that
\[
\Delta_{V_1}c(\emptyset) = c(\{V_1\})-c(\emptyset) = \jnum{n5.delta.small},
\qquad
\Delta_{V_1}c(\{V_2\}) = c(\{V_1,V_2\})-c(\{V_2\}) = \jnum{n5.delta.big}.
\]
Submodularity of a cost requires $\Delta_x c(A)\ge\Delta_x c(B)$ for $A\subseteq B$; here
$\jnum{n5.delta.small} \ge \jnum{n5.delta.big}$ is false, so $c$ is supermodular on this
pair. The mechanism is transparent: under $\pi_0$ both vehicles charge in the expensive band;
a single coordinating vehicle shifts its session into the cheap window and saves; when
\emph{both} coordinate, only one of them can still have the window, so the second gains
nothing.

For the induced savings game,
$v(\{V_1\})=v(\{V_2\})=v(\{V_1,V_2\})=\jnum{n5.v}$, so a core allocation would require
$x_1+x_2=\jnum{n5.v}$ with $x_1\ge\jnum{n5.v}$ and $x_2\ge\jnum{n5.v}$ --- impossible for a
positive value. The core is empty and the least-core radius is $\jnum{n5.epsilon}$. The
Shapley value $(\jnum{n5.v.half},\jnum{n5.v.half})$ is not in it.
\end{proof}

The construction isolates the mechanism; the mechanism itself is general. \emph{Any} scarce
shared resource that saturates converts diminishing returns into congestion, and congestion
is supermodular. The control experiment confirms the causal direction: duplicating the
charging station in the same instance restores the second vehicle's marginal saving from
zero to $\jnum{n5.gain.relaxed}$ (Section~\ref{sec:res-coop}).

This is not a defect to engineer around. It is a structural property of the problem, and it
decides the architecture of the analysis: the headline stability result must be one that is
\emph{always} available.

\begin{proposition}[Conditional convexity]
\label{prop:p5}
On the congestion-free sub-class --- instances in which no shared resource saturates over the
horizon --- $c$ is submodular, hence $v$ is convex, hence the core is non-empty and
$\varphi(v)\in \mathcal{C}(v)$~\citep{shapley1971cores}.
\end{proposition}

We state Proposition~\ref{prop:p5} as a conditional proposition and are explicit that it is
not established unconditionally here. Two obligations stand between the submodularity of the
stage objective and that of $c$, and only the first is discharged in this paper. The
congestion-free condition must be stated verifiably from a simulation trace, which it is
(saturation is measured directly as resource utilisation). But the descent of submodularity
from the stage objective to $c$ through the min-envelope does not hold automatically: a
pointwise minimum of submodular functions need not be submodular, and the two functions have
opposite monotonicity --- $\widehat{\Phi}$ is non-decreasing in committed (player, action)
pairs while $c$ is non-increasing in coalition membership. That is a genuine proof step, not a
restriction, and we do not claim to have closed it. Section~\ref{sec:res-coop} therefore
treats Proposition~\ref{prop:p5} as an \emph{empirical} claim certified per instance, and
nothing else in the paper depends on it.

\subsection{Certified stability}

\begin{theorem}[Least-core certification]
\label{thm:leastcore}
For every instance the programme
\begin{equation}
\max\ \varepsilon \quad\text{s.t.}\quad
\sum_{i\in S}x_i \ \ge\ v(S)+\varepsilon \ \ \forall\, \emptyset\neq S\subset N,
\qquad \sum_{i\in N}x_i = v(N)
\label{eq:leastcore}
\end{equation}
is feasible and bounded, and its optimum $\varepsilon^\star$ satisfies: $\varepsilon^\star\ge0$
if and only if the core is non-empty, and $\varepsilon^\star<0$ is the least-core radius ---
the largest shortfall any coalition suffers under the most stable allocation.
\end{theorem}

The direction of~\eqref{eq:leastcore} is the thing to get right. $v$ is a savings (profit)
game, so an allocation is blocked by a coalition that could secure \emph{more} on its own
than it is being given, which is the $\ge$ above. The cost-game form
$\sum_{i\in S}x_i \le v(S)+\varepsilon$, with ``$\varepsilon\le0$ certifies membership'',
tests the wrong condition with an inverted sign. The released test suite pins the direction
against a hand-computed three-player example and against the instance of
Proposition~\ref{prop:n5}, whose core must come out empty.

For $n=|\mathcal{M}|+|\mathcal{V}|$ small enough --- and on this benchmark $n=6$, so
$2^{6}=64$ coalition evaluations --- the Shapley value, the least core and the nucleolus are
all computed \emph{exactly}, with no sampling error. A stratified, antithetic
sampler~\citep{castro2009polynomial} with reported confidence intervals is available for
larger games, and constraint generation replaces full enumeration in~\eqref{eq:leastcore}.

\begin{proposition}[Monotonicity]
\label{prop:t6}
$v(S)\le v(S')$ for $S\subseteq S'$ holds by construction for the exact characteristic
function, since a larger coalition can replicate the smaller one's policy.
\end{proposition}

Proposition~\ref{prop:t6} is stated for the exact $c$. It need \emph{not} hold for the
single-sweep surrogate actually used, because a greedy sweep with more participants is not
guaranteed to dominate one with fewer. We therefore measure the surrogate's violation rate
rather than assuming it away; it is \jnum{e3.mono.violation.mean}\,\% of sampled pairs.

\subsection{Layer 1: decentralised coalition-structure generation}
\label{sec:layer1}

The grand coalition requires all-to-all coordination, which contradicts the decentralisation
thesis and does not scale. Layer~1 finds a partition $CS=\{S_1,\dots,S_K\}$ of $N$.

Agent $i$ prefers coalition $S$ to $S'$ iff $\varphi_i(v|_S) > \varphi_i(v|_{S'})$, i.e.\
hedonic preferences over the within-coalition Shapley allocation, where
$v|_S(T)=c(\emptyset)-c(T)$ for $T\subseteq S$. Starting from singletons, merge--split
\citep{apt2009stable, saad2009coalitional} accepts any union that is Pareto-improving for all
its members and any split of a coalition that is Pareto-improving for all of its; it
terminates because the Pareto order over partitions is acyclic, and its fixed point is
$\mathbb{D}_{hp}$-stable: no group can profitably merge and no coalition can profitably split.
A cap $S_{\max}$ on coalition size bounds the negotiation neighbourhood, and the number of
bilateral negotiations actually conducted is reported as the communication cost.

Because congestion can make the grand coalition non-convex, coalition formation here is not
merely a tractability device: when a shared resource saturates, partitioning can be genuinely
\emph{better} than the grand coalition, and Section~\ref{sec:res-coop} tests exactly that.

One distinction matters for reporting. $\sum_k v(S_k)$ is an accounting figure --- it adds
savings each measured against a $\pi_0$ complement --- whereas the value a partition actually
delivers must be measured by running all blocks simultaneously, each coordinating internally
only. We report the measured quantity.

\subsection{Layer 3: bargaining over the trade-off weight}
\label{sec:layer3}

Partition the agents into two interest blocs: production $\mathcal{M}$, whose payoff decreases
in $C_{\max}$, and logistics/energy $\mathcal{V}$, whose payoff decreases in
$E_{\mathrm{cost}}$. Let the attainable set be the Pareto front of equilibrium outcomes swept
over $\omega$.

The disagreement point requires care. Taken literally, ``the payoffs at the $M1$ equilibrium''
is a point \emph{on} the front, so the bargaining set is empty and neither solution is
defined. The equilibrium payoffs that actually bound the negotiation are the two extremes:
production's worst outcome is the makespan it suffers when logistics dictates
($\omega=0$), and logistics' worst is the bill it pays when production dictates
($\omega=1$). That pair is the nadir of the equilibrium front, and it is the disagreement
point $d$ used here.

The Nash bargaining solution~\citep{nash1950bargaining} maximises
$(u_{\mathcal{M}}-d_{\mathcal{M}})(u_{\mathcal{V}}-d_{\mathcal{V}})$ over the front; the
Kalai--Smorodinsky solution~\citep{kalai1975other} equalises the two normalised gains. The
supporting hyperplane of the front at the selected point yields an interval of weights
$[\omega^{-},\omega^{+}]$ for which that point minimises the normalised scalarisation, and
$\omega^{\mathrm{NBS}}$ is reported as the midpoint together with the interval. An empty
interval would mean the selected point lies in a non-convex pocket that no weighted sum can
reach, which is itself worth reporting.

\begin{proposition}[Scale invariance]
\label{prop:scale}
Both solutions are invariant under independent positive affine rescaling of each bloc's
utility. Hence Layer~3 returns the same operating point under any choice of units or
normalisation.
\end{proposition}

\begin{proof}
Let $u'_j=\alpha_j u_j+\beta_j$ with $\alpha_j>0$, so $d'_j=\alpha_j d_j+\beta_j$. Then
$\prod_j (u'_j-d'_j)=\big(\prod_j\alpha_j\big)\prod_j(u_j-d_j)$, a positive constant multiple,
so the maximiser is unchanged. For Kalai--Smorodinsky, the ideal point transforms the same
way, so each ratio $(u'_j-d'_j)/(u^{\mathrm{ideal}\prime}_j-d'_j)$ equals its untransformed
counterpart, and the equalising point is unchanged.
\end{proof}

Proposition~\ref{prop:scale} is why Layer~3, alone among the scalarised machinery of this
paper, required no repair for commensurability: it never adds the two objectives together.
It is also the strongest argument for deriving $\omega$ rather than choosing it.

\paragraph{Managerial variant.} Where a single unit is wanted,
$\Phi_{\text{\euro}} = c_{\mathrm{time}}C_{\max} + E_{\mathrm{cost}}$ has no weight at all,
and $c_{\mathrm{time}}$ --- the opportunity cost of shop time --- is a quantity a plant knows.
Layer~3 then derives $c_{\mathrm{time}}$ instead of $\omega$, and the break-even value at
which the cost-optimal schedule flips from off-peak-seeking to deadline-seeking is the number
to report.

\section{Experimental design}
\label{sec:exp}

\subsection{Instances}
The evaluation uses the \jnum{bench.instances} instances of the Bilge--Ulusoy
family~\citep{bilge1995}: ten job sets crossed with four travel-time layouts, four machines,
and two vehicles. Job sets, processing times and layouts are transcribed verbatim. The fleet
size is nowhere stated in the source of the tariff profiles and is recovered from the horizon
formula~\eqref{eq:horizon}: job set~1 with layout~1 gives
$2(176/4+128/|\mathcal{V}|)=216.0$, hence $|\mathcal{V}|=2$. The recovery is verified rather
than assumed --- the regenerated horizon reproduces the published value for all
\jnum{bench.instances} instances.

The tariff profile is \emph{regenerated} from~\eqref{eq:horizon} rather than copied, because
seven rows of the published table have a sixth period ending after their own stated horizon,
which contradicts the feasibility condition the same source imposes on individual actions.
The regeneration reproduces the published boundaries of periods one to five on all
\jnum{bench.instances} instances --- which is the transcription check --- and, unlike the
published table, always terminates the last period at $H$.

\subsection{Methods}
\begin{description}
\item[$\pi_0$] the reference completion policy: shortest buffered operation, earliest-delivery
transport, charge only when no transport is battery-feasible, no deferral. This is the policy
against which marginal contributions and coalition costs are measured, so it is reported in
its own right.
\item[\textsc{dr-*}] a family of \jnum{exp.rules.count} composite dispatching rules --- five
machine rules $\times$ four vehicle rules $\times$ three charging policies $\times$ two
deferral policies. We report a fixed member (\textsc{dr-spt}), the energy-aware member
(\textsc{dr-tou}), and \textsc{dr-oracle}: the best of all \jnum{exp.rules.count} chosen
\emph{per instance} with hindsight. \textsc{dr-oracle} is not implementable; it is included
because comparing against a single hand-picked rule would be a strawman.
\item[$M1a$, $M1b$] best response and log-linear learning.
\item[$M2$] the grand coalition under Definition~\ref{def:c}, and the merge--split partition.
\item[\textsc{ex-cp}] the exact reference of Section~\ref{sec:exact}.
\end{description}

\subsection{Protocol and statistics}
All stochastic components take an explicit generator; no global seed is read, and event ties
break on a documented total order, so a given seed reproduces a run exactly. The weight grid
is $\omega\in\{0.25,0.5,0.75\}$ on the full benchmark, with a finer grid for the bargaining
front. The comparison follows the standard protocol for multiple methods over multiple
problems~\citep{demsar2006statistical}: a Friedman omnibus test across methods over
instances, and where it is significant, pairwise Wilcoxon signed-rank tests with
Holm--Bonferroni correction, each accompanied by a Vargha--Delaney $A_{12}$ effect
size~\citep{vargha2000critique}. No mean is reported without a dispersion measure, and no
comparison is described as significant without the corrected test.

\subsection{Metrics}
Beyond $\Phi$, $C_{\max}$ and $E_{\mathrm{cost}}$ we report: the deadline-feasibility rate and
the deadlock rate; the least-core radius $\varepsilon^\star$ and its ratio to $v(N)$; the
fraction of instances whose core is non-empty and whose Shapley value is core-stable; the
submodularity violation rate; the Gini coefficient of the allocation; coalition sizes and the
number of bilateral messages; charger utilisation; and per-decision wall-clock.

\section{Results}
\label{sec:results}

\subsection{Commensurability: what the raw objective was measuring}
\label{sec:res-e13}
Table-free but decisive. Across the benchmark the raw makespan and energy terms differ by a
factor of \jnum{e13.ratio.min}--\jnum{e13.ratio.max} (mean \jnum{e13.ratio.mean}). At
$\omega=0.5$ the raw energy term therefore contributes a mean of
\jnum{e13.raw.energy.share.mean}\,\% of the objective, at most
\jnum{e13.raw.energy.share.max}\,\% on any instance, and the two terms balance only near
$\omega^\star=\jnum{e13.omega.star.mean}$. A weighted sum on the raw quantities is a makespan
objective for every weight above roughly $0.02$, and any sweep over $\omega\in[0,1]$ would be
flat. After the normalisation of~\eqref{eq:phi} the energy term contributes
\jnum{e13.norm.energy.share.mean}\,\% at the same weight, and the weight becomes an actual
control --- as Section~\ref{sec:res-e1} shows by moving it.

\subsection{Where the exact model stops closing}
\label{sec:res-exact}
Under a per-instance budget of \jnum{e2.budget.seconds}~s, \textsc{ex-cp} closes the
classical makespan variant on \jnum{e2.makespan.closed.rate}\,\% of the tested instances, in
at most \jnum{e2.makespan.seconds.max}~s, and reproduces the original solver's optimum
exactly. Adding the ToU cost and the deadline drops the closure rate to
\jnum{e2.tou.closed.rate}\,\%; adding charging, charger capacity and the state-of-charge
recursion drops it to \jnum{e2.full.closed.rate}\,\%, with the model growing from
\jnum{e2.makespan.vars.mean} to \jnum{e2.full.vars.mean} variables.

The mechanism is worth naming, because it is not simply ``a bigger model''. Minimising
makespan alone leaves a large plateau of equally optimal schedules and CP-SAT exploits it;
pricing energy against the tariff makes every member of that plateau distinguishable, and the
state-of-charge recursion couples the vehicle routes that the classical model could treat
almost independently. The bounds \textsc{ex-cp} returns on the full model within this budget
are too weak to support a price-of-anarchy estimate, and we do not report one. This is the
concrete scalability wall that motivates decentralisation, and it is reported as evidence;
the honest consequence is that Theorem~\ref{thm:t3} stands as proved while the empirical
price of anarchy remains open (Section~\ref{sec:limitations}).

\subsection{Method comparison}
\label{sec:res-e1}
Table~\ref{tab:e1} reports the main comparison at $\omega=0.5$. The Friedman test across the
seven methods over \jnum{bench.instances} instances is significant
($\chi^2=\jnum{e1.friedman.stat}$, $p=\jnum{e1.friedman.p}$), with a Nemenyi critical
difference of \jnum{e1.friedman.cd} mean ranks. The diagram separates the two game methods
and the rule oracle, as a group, from everything else; it does not separate them from each
other.

$M2$ in its grand-coalition form has the best mean rank (\jnum{e1.rank.m2grand}), reducing
$\Phi$ by \jnum{e1.m2grand.gain.mean}\,\% $\pm$ \jnum{e1.m2grand.gain.std}\,\% against
$\pi_0$ and beating the hindsight-selected \textsc{dr-oracle} on
\jnum{e1.m2grand.winrate.vs.oracle}\,\% of instances. $M1a$ follows at rank
\jnum{e1.rank.m1a} with a \jnum{e1.m1a.gain.mean}\,\% $\pm$ \jnum{e1.m1a.gain.std}\,\%
reduction.

The tests must be read carefully, and they do not all say the same thing. Against $\pi_0$,
$M2$ wins decisively: Wilcoxon with Holm correction gives
$p=\jnum{e1.wilcoxon.m2vspi0.p}$ with $A_{12}=\jnum{e1.wilcoxon.m2vspi0.a12}$, and the rank
gap exceeds the critical difference. Against \textsc{dr-oracle} the same test gives
$p=\jnum{e1.wilcoxon.m2vsdroracle.p}$ ($A_{12}=\jnum{e1.wilcoxon.m2vsdroracle.a12}$), but
the rank gap is smaller than the Nemenyi critical difference, so the conservative
all-pairs post-hoc does \emph{not} separate them. \textbf{And $M2$ is not significantly
better than $M1a$}: $p=\jnum{e1.wilcoxon.m2vsm1a.p}$,
$A_{12}=\jnum{e1.wilcoxon.m2vsm1a.a12}$, a rank gap far inside the critical difference. On
this benchmark the cooperative layer's contribution is \emph{stability analysis and
allocation}, not a significant improvement in objective value over the non-cooperative
mechanism; that is what the data support and it is what we claim.

Three observations deserve emphasis because they are not the expected ones.

First, \textbf{$M1a$ does not dominate the rule oracle}: it wins on
\jnum{e1.m1a.winrate.vs.oracle}\,\% of instances. The best of
\jnum{exp.rules.count} composite rules chosen with hindsight is a strong comparator, and we
report the win rate rather than claiming dominance. Against the implementable comparators ---
$\pi_0$ and any fixed rule --- the mechanism wins throughout.

Second, \textbf{the energy-greedy dispatching rule is not a good scheduler}. \textsc{dr-tou}
has the \emph{worst} mean rank of all methods (\jnum{e1.rank.drtou}) despite optimising the
right quantity locally. It achieves what it sets out to: its mean energy bill,
\jnum{e1.drtou.ecost.mean}~\euro, is the lowest of any method, below even $M2$'s
\jnum{e1.m2grand.ecost.mean}~\euro. It pays for that with the deadline --- a mean makespan of
\jnum{e1.drtou.cmax.mean}~min against \jnum{e1.m2grand.cmax.mean}~min, and a
deadline-feasibility rate of \jnum{e1.drtou.deadline.rate}\,\% against
\jnum{e1.m1a.deadline.rate}\,\% for $M1a$. Energy awareness without a feasibility safeguard
is worse than no energy awareness at all, which is precisely the risk that motivates the
remaining-work bound of Section~\ref{sec:deadline}.

Third, \textbf{best response converges immediately}. $M1a$ needs
\jnum{e1.m1a.iterations.mean} round-robin iterations per event on average over
\jnum{e1.m1a.stages.mean} events, and the largest $\varepsilon$-Nash certificate observed
anywhere in the benchmark is \jnum{e1.m1a.epsne.max} --- the committed profiles are exact
stage equilibria, as Theorem~\ref{thm:t1} requires. The iteration cap never binds.

The deadlock rates make the second point sharper. The safeguard of
Section~\ref{sec:deadline} never fires for $\pi_0$, for the game methods or for
\textsc{dr-spt} (\jnum{e1.m1a.fallback.rate}\,\% for $M1a$), so their deadline results are
not an artefact of it. It fires on \jnum{e1.drtou.fallback.rate}\,\% of \textsc{dr-tou}'s
episodes --- the energy-greedy rule drives itself into states with no feasible completion on
nearly three quarters of the benchmark, and even with the safeguard it still finishes late on
$100-\jnum{e1.drtou.deadline.rate}$\,\% of instances. Risk R13 is not hypothetical.

\begin{table}[t]
\centering
\small
\caption{Main comparison on the benchmark at $\omega=0.5$. Lower $\Phi$ is better; ranks are
Friedman mean ranks. \textsc{dr-oracle} is the best of \jnum{exp.rules.count} composite rules
selected per instance with hindsight and is not implementable.}
\label{tab:e1}
\begin{tabular}{@{}lrrrrrr@{}}
\toprule
Method & $\Phi$ mean & $\Phi$ s.d. & $C_{\max}$ & $E_{\mathrm{cost}}$ & Deadline & Rank \\
       &             &             & (min)      & (\euro)             & met (\%) &      \\
\midrule
$\pi_0$              & \jnum{e1.pi0.phi.mean}      & \jnum{e1.pi0.phi.std}      & \jnum{e1.pi0.cmax.mean}      & \jnum{e1.pi0.ecost.mean}      & \jnum{e1.pi0.deadline.rate}      & \jnum{e1.rank.pi0} \\
\textsc{dr-spt}      & \jnum{e1.drspt.phi.mean}    & \jnum{e1.drspt.phi.std}    & \jnum{e1.drspt.cmax.mean}    & \jnum{e1.drspt.ecost.mean}    & \jnum{e1.drspt.deadline.rate}    & \jnum{e1.rank.drspt} \\
\textsc{dr-tou}      & \jnum{e1.drtou.phi.mean}    & \jnum{e1.drtou.phi.std}    & \jnum{e1.drtou.cmax.mean}    & \jnum{e1.drtou.ecost.mean}    & \jnum{e1.drtou.deadline.rate}    & \jnum{e1.rank.drtou} \\
\textsc{dr-oracle}   & \jnum{e1.droracle.phi.mean} & \jnum{e1.droracle.phi.std} & \jnum{e1.droracle.cmax.mean} & \jnum{e1.droracle.ecost.mean} & \jnum{e1.droracle.deadline.rate} & \jnum{e1.rank.droracle} \\
$M1b$                & \jnum{e1.m1b.phi.mean}      & \jnum{e1.m1b.phi.std}      & \jnum{e1.m1b.cmax.mean}      & \jnum{e1.m1b.ecost.mean}      & \jnum{e1.m1b.deadline.rate}      & \jnum{e1.rank.m1b} \\
$M1a$                & \jnum{e1.m1a.phi.mean}      & \jnum{e1.m1a.phi.std}      & \jnum{e1.m1a.cmax.mean}      & \jnum{e1.m1a.ecost.mean}      & \jnum{e1.m1a.deadline.rate}      & \jnum{e1.rank.m1a} \\
$M2$ (grand)         & \jnum{e1.m2grand.phi.mean}  & \jnum{e1.m2grand.phi.std}  & \jnum{e1.m2grand.cmax.mean}  & \jnum{e1.m2grand.ecost.mean}  & \jnum{e1.m2grand.deadline.rate}  & \jnum{e1.rank.m2grand} \\
\bottomrule
\end{tabular}
\end{table}

\subsection{The cooperative game}
\label{sec:res-coop}

\paragraph{Savings.} Coordinating the whole shop reduces $\Phi$ by
\jnum{e3.saving.mean}\,\% $\pm$ \jnum{e3.saving.std}\,\% relative to $\pi_0$, up to
\jnum{e3.saving.max}\,\% on the most responsive instance, giving a mean grand-coalition value
$v(N)=\jnum{e3.v.mean}$.

The dispersion hides something that should be said plainly. On
\jnum{e3.negative.saving.count} of \jnum{bench.instances} instances the grand coalition is
\emph{worse} than $\pi_0$ ($v(N)\le0$, worst case \jnum{e3.saving.min}\,\%). This is a
property of the surrogate, not of coordination: $c(S)$ is one greedy best-response sweep, and
a greedy sweep with more participants is not guaranteed to dominate one with fewer --- the
same reason Proposition~\ref{prop:t6} is claimed for the exact $c$ and only measured for the
surrogate. Those instances are excluded from the ratio statistics below, where
$\varepsilon^\star/v(N)$ would change sign for a reason unrelated to stability; the remaining
\jnum{e3.eps.ratio.n} carry them.

\paragraph{The core is empty, throughout.} This is the central empirical finding.
$\varepsilon^\star<0$ on \jnum{e3.core.empty.count} of \jnum{bench.instances} instances
(\jnum{e3.core.empty.rate}\,\%); the core is non-empty on
\jnum{e3.core.nonempty.rate}\,\% of them, and the Shapley value is core-stable on
\jnum{e3.shapley.incore.rate}\,\%. The mean least-core radius is
$\varepsilon^\star/v(N)=\jnum{e3.eps.ratio.mean}$ (median \jnum{e3.eps.ratio.median}, worst
\jnum{e3.eps.ratio.worst}), i.e.\ the most stable allocation still leaves some coalition
short by roughly half of what it could secure alone. The submodularity sampling test agrees:
the diminishing-returns inequality is violated on \jnum{e3.submod.violation.mean}\,\% of
sampled triples on average (range \jnum{e3.submod.violation.min}--%
\jnum{e3.submod.violation.max}\,\%).

Two consequences follow, and neither is negotiable. Any paper in this problem class that
reports a Shapley allocation as ``stable'' because the game ``should be'' convex is reporting
something false. And the least-core radius is not a fallback but the correct headline: it is
computable on every instance, it is signed correctly, and it quantifies exactly how much
instability an allocation carries.

\paragraph{Which resource saturates.} The constructed instance of
Proposition~\ref{prop:n5} makes the charger the culprit, and there the causal control is
clean: with one station the second coordinating vehicle gains
$\jnum{n5.gain.congested}$, with two it gains $\jnum{n5.gain.relaxed}$, and the least-core
radius improves by \jnum{n5.epsilon.improvement}\,\%. On the benchmark, however, the charger
is \emph{not} the binding resource: its utilisation is only
\jnum{e11.charger1.utilisation}\,\% under the reference policy, and duplicating it moves the
mean least-core radius from \jnum{e11.charger1.eps.mean} to \jnum{e11.charger2.eps.mean} and
the submodularity violation rate from \jnum{e11.charger1.submod.violation}\,\% to
\jnum{e11.charger2.submod.violation}\,\% --- that is, not at all.

The obstruction is therefore \emph{not charger-specific}. What saturates on these instances is
the bottleneck machine and the two-vehicle fleet, and they behave exactly as
Proposition~\ref{prop:n5} predicts a saturating shared resource must. The generalisation
matters more than the special case: adding charging capacity will not buy stability in a shop
whose machines are the constraint.

\paragraph{Allocation.} Vehicles capture \jnum{e3.robot.share.mean}\,\% $\pm$
\jnum{e3.robot.share.std}\,\% of the total savings --- a large dispersion, driven by the few
instances on which the surrogate makes coordination lose --- (mean Shapley payoff
\jnum{e3.robot.payoff.mean} against \jnum{e3.machine.payoff.mean} for machines), with a Gini
coefficient of \jnum{e3.gini.shapley.mean}. The asymmetry is structural rather than a
modelling artefact: vehicles hold both levers --- they choose when to charge, and their
routing decides when parts arrive --- whereas a machine controls only the deferral of an
operation whose position in the sequence is already fixed by the vehicles.

\paragraph{Cost of certification.} Exact certification required
\jnum{e3.calls.per.instance} coalition evaluations per instance, \jnum{e3.calls.total} in
total, at \jnum{e3.ms.per.call}~ms each for a total of \jnum{e3.seconds.total}~s. No sampling
was needed at $n=6$; the estimator and its confidence intervals are implemented for larger
games and are not exercised here. The fraction of coalitions that miss the deadline is
\jnum{e3.infeasible.coalition.rate}\,\%, so the penalised characteristic function is
essentially inactive on this benchmark --- the coalition structure is organised around the
tariff, not around the deadline.

\subsection{Coalition structures}
\label{sec:res-cs}
Merge--split converges to a $\mathbb{D}_{hp}$-stable partition on
\jnum{e4.dhp.stable.rate}\,\% of instances, with \jnum{e4.blocks.mean} blocks on average, a mean largest block of
\jnum{e4.maxblock.mean} agents (\jnum{e4.maxblock.max} at most), and
\jnum{e4.messages.mean} bilateral negotiations. That partition realises \jnum{e4.value.ratio.mean}\,\% of the grand-coalition
savings on average and \jnum{e4.value.ratio.median}\,\% at the median (worst case
\jnum{e4.value.ratio.min}\,\%).

This is the quantitative content of the decentralisation claim: almost all of the value of
full coordination is available through \emph{pairwise} agreements, at a communication cost
that does not grow with the shop. And on \jnum{e4.partition.beats.grand.count} instances the partition is strictly
\emph{better} than the grand coalition --- impossible for a convex game, and a direct
empirical corollary of Proposition~\ref{prop:n5}: when a shared resource saturates, adding a
member to a coalition can hurt.

\subsection{Bargaining-derived weight}
\label{sec:res-barg}
Over the equilibrium front, Nash bargaining selects
$\omega^{\mathrm{NBS}}=\jnum{e5.omega.nbs.mean}\pm\jnum{e5.omega.nbs.std}$ (median
\jnum{e5.omega.nbs.median}); Kalai--Smorodinsky selects \jnum{e5.omega.ks.mean}, and the two
agree on the operating point on \jnum{e5.nbs.ks.agree.rate}\,\% of instances. Fronts contain
\jnum{e5.front.size.mean} distinct non-dominated points on average.

The number is the result. A shop that splits the difference between throughput and energy
would set $\omega=0.5$; bargaining between the two blocs, with the non-cooperative extremes
as the threat point, selects roughly \jnum{e5.omega.nbs.mean} --- markedly more
energy-weighted. Under the monetised variant the same front gives a break-even opportunity
cost of shop time of \jnum{e5.breakeven.ctime.mean}~\euro/min, i.e.\
\jnum{e5.breakeven.ctime.hour}~\euro/h: below that value a plant should chase off-peak
windows, above it, throughput. That is a quantity a plant can argue about, and it is derived
rather than assumed.

\subsection{Ablations}
\label{sec:res-lambda}
\textbf{The horizon parameter $\lambda$.} At the published $\lambda=2$ every instance is
deadline-admissible (\jnum{e11.lambda20.admissible}\,\% with margin $\ge1.05$) but the
deadline is \emph{binding} --- margin $\le 2$ --- on only \jnum{e11.lambda20.binding}\,\% of
them, at a mean margin of \jnum{e11.lambda20.margin.mean}. Tightening to $\lambda=1.5$ makes
it binding on \jnum{e11.lambda15.binding}\,\% while keeping
\jnum{e11.lambda15.admissible}\,\% admissible; at $\lambda=1.25$ the deadline binds
everywhere (\jnum{e11.lambda125.binding}\,\%) at the cost of
$100-\jnum{e11.lambda125.admissible}$\,\% of instances becoming inadmissible.

So $\lambda$ is a feasibility parameter, not a presentational one, and at its published value
the deadline is present but slack. We report the results at $\lambda=2$ for comparability
with the published profiles and note that the deadline's role here is to bound the deferral
lever rather than to be the active constraint.

\textbf{The deadline penalty $M$.} Varying $M$ over $\{2,5,10,50\}$ moves the Gini
coefficient of the Shapley allocation by \jnum{e11.penalty.gini.spread.max} and the
least-core radius by \jnum{e11.penalty.eps.spread.max} --- that is, not at all, on any
instance. The reason is worth stating rather than presenting as a robustness result: the
infeasible-coalition rate is \jnum{e3.infeasible.coalition.rate}\,\%, so on this benchmark
the penalty term is never active and $M$ has nothing to scale. The ablation establishes
that the allocations are not artefacts of the penalty, which is what it exists to do; it
does not establish that $M$ would be harmless on a benchmark where the deadline binds
harder, and we do not claim that.

\textbf{Charger capacity $K_{CH}$.} Reported in Section~\ref{sec:res-coop}: on this benchmark
it changes nothing, because the charger does not saturate.

\subsection{Limitations}
\label{sec:limitations}
Four limitations are load-bearing and we state them here rather than in the conclusion.

\emph{(i) No empirical price of anarchy.} \textsc{ex-cp} does not close the full model within
the budget and its bounds are too weak to normalise against
(Section~\ref{sec:res-exact}). Theorem~\ref{thm:t3} is proved and tested;
the price of anarchy is measured nowhere in this paper and no surrogate for it is offered.

\emph{(ii) $c(S)$ is a surrogate.} The characteristic function is evaluated by a single
best-response sweep, not by exact optimisation of the coalition's problem. It is
deterministic, and every solution concept computed from it is exact \emph{given} it, but the
gap between this surrogate and the true $c$ is not characterised --- for the same reason as
(i). Two measurable symptoms are reported rather than smoothed over: the surrogate violates
the monotonicity that the exact $c$ satisfies by construction, on
\jnum{e3.mono.violation.mean}\,\% of sampled pairs; and on
\jnum{e3.negative.saving.count} instances it makes the grand coalition lose to $\pi_0$
outright. Both would be impossible for the exact characteristic function.

\emph{(iii) Proposition~\ref{prop:p5} is conditional and unproved.} The min-envelope step is
open (Section~\ref{sec:n5}). Nothing in the paper's conclusions rests on it.

\emph{(iv) Scale.} The benchmark is the classical four-machine family, which is what makes
exact cooperative certification affordable, but it cannot by itself evidence behaviour at
industrial scale. A seeded generator for larger instances is included in the released
artefact; the certification results here are stated for $n=6$ and are not extrapolated.

\section{Practical and managerial insights}
\label{sec:insights}

\paragraph{Energy awareness needs a feasibility safeguard, or it backfires.}
The single most actionable result is negative. A dispatching rule that greedily shifts load
into off-peak bands --- exactly what a plant would implement first --- produced the
\emph{worst} outcomes of any method tested (Section~\ref{sec:res-e1}), because deferral
consumes deadline slack that cannot be recovered. Deferral must be paired with a check on
remaining work against remaining time. In the mechanism proposed here that check is intrinsic:
an action that would make the deadline unreachable is not in the agent's choice set.

\paragraph{Fleet decisions dominate machine decisions.}
Vehicles capture \jnum{e3.robot.share.mean}\,\% of the coordination value. For a plant
choosing where to invest in coordination technology, the ordering is clear: automating the
vehicles' charge-timing and routing decisions returns several times what automating machine
deferral does, because a machine's deferral options are bounded by an arrival sequence the
vehicles have already fixed.

\paragraph{Pairwise agreements are enough.}
The stable coalition structures found here contain at most
\jnum{e4.maxblock.max} agents, \jnum{e4.maxblock.mean} on average, and require
\jnum{e4.messages.mean} bilateral negotiations, yet
capture \jnum{e4.value.ratio.mean}\,\% of the value of coordinating everything. A shop does
not need a central scheduler or a shop-wide consensus protocol; it needs a way for a machine
and the vehicle that feeds it to agree.

\paragraph{Do not promise a stable cost allocation.}
The core is empty on \jnum{e3.core.empty.rate}\,\% of instances, so \emph{no} allocation of
the joint savings is immune to renegotiation, the Shapley value included. Internal transfer
pricing between a production and a logistics cost centre should therefore be set from the
least core --- the allocation minimising the worst coalition's grievance --- and the residual
grievance \jnum{e3.eps.ratio.mean} of total savings should be disclosed to both parties
rather than papered over. The practical form of this is a periodic review clause, not a fixed
split.

\paragraph{The trade-off weight has a defensible value.}
Rather than assuming $\omega=0.5$, bargaining between the two blocs selects
\jnum{e5.omega.nbs.mean}. Equivalently, in monetary terms, off-peak seeking pays whenever the
opportunity cost of shop time is below \jnum{e5.breakeven.ctime.hour}~\euro/h. A plant that
knows its own contribution margin per hour can read its policy directly off that number, and
the answer is invariant to the units in which either objective is measured
(Proposition~\ref{prop:scale}).

\paragraph{Adding chargers is not a stability lever.}
On this benchmark the charging station is idle \jnum{e11.charger1.utilisation}\,\% of the
time and duplicating it changes neither the least-core radius nor the submodularity violation
rate. Where the machines are the bottleneck, charging capacity is not what is scarce, and
capital spent there buys nothing in coordination terms.

\section{Conclusion}
\label{sec:conc}

We have developed a decentralised mechanism for scheduling production machines and
battery-constrained autonomous vehicles under time-of-use tariffs, with the planning horizon
treated as a hard deadline and both energy streams priced against the tariff. On the
non-cooperative side, each stage game with marginal-contribution utilities is an exact
potential game whose price of stability is one, so sequential best response converges to a
pure equilibrium --- and we are explicit that this is a stage-level result which does not
extend to trajectory optimality, and that the frequently quoted $\tfrac12$-suboptimality
bound cannot be transposed from a maximisation of submodular welfare to a cost objective.

The cooperative contribution is where the surprise lies. The natural hypothesis --- that
coordination has diminishing returns, so the coalition cost is submodular, the savings game
convex, and the core non-empty with the Shapley value inside --- is false for this problem.
Congestion in any saturating shared resource makes the cost supermodular, and we exhibit an
instance whose core is provably empty and confirm that the obstruction is present on every
benchmark instance. Stability must therefore be \emph{certified} rather than asserted, and the
least-core radius is the right instrument because it is always computable and correctly
signed. Reporting a certified radius is a weaker claim than a core-existence theorem, and a
truer one.

Around that result, decentralised merge--split coalition formation recovers almost all of the
grand-coalition value through pairwise agreements, and Nash bargaining between the production
and logistics blocs supplies the trade-off weight that scalarised bi-objective scheduling
normally leaves to the author.

Three directions follow directly. The exact reference must be strengthened before an
empirical price of anarchy can be reported; a stronger relaxation or a decomposition of the
state-of-charge recursion is the natural target. The min-envelope step of
Proposition~\ref{prop:p5} remains open and would turn a conditional proposition into a
theorem. And the mechanism's per-decision cost --- a few hundred milliseconds at this scale
--- suggests replacing the counterfactual rollout with a learned value estimate, which is the
subject of a companion study: because the marginal-contribution utility that makes this game
a potential game is definitionally the difference reward used for multi-agent credit
assignment, a learner trained on that signal performs policy gradient on the same potential
the mechanism descends.

\section*{Reproducibility}
The simulator, all solution concepts, the benchmark files with regenerated tariff profiles,
the experiment runners and the test suite are released. Every numerical value in this paper
resolves to a key in a machine-readable registry produced by a recorded run; no numeral in a
results sentence is typed by hand. The test suite contains, for each formal result, a test
that fails if the result is false --- including a regression test pinning the counterexample
of Proposition~\ref{prop:n5}, whose core must come out empty --- and a regression test
asserting that under a flat tariff with no deadline and no battery the exact model reproduces
the original solver's makespan exactly.

\section*{CRediT authorship contribution statement}
\textbf{Kader Sanogo}: Conceptualization, Methodology, Software, Formal analysis,
Investigation, Writing -- original draft. \textbf{Malek Masmoudi}: Conceptualization,
Supervision, Writing -- review \& editing. \textbf{Abdelkader Mekhalef Benhafssa}:
Methodology, Validation, Writing -- review \& editing. \textbf{M'hammed Sahnoun}:
Supervision, Funding acquisition, Writing -- review \& editing.

\section*{Declaration of competing interest}
The authors declare that they have no known competing financial interests or personal
relationships that could have appeared to influence the work reported in this paper.

\appendix

\section{Notation}
\label{app:notation}
\begin{table}[h!]
\centering
\small
\caption{Notation. Symbol collisions present in earlier drafts of this work are resolved
here: the semi-Markov tuple is written $\mathfrak{M}$ rather than $\mathcal{M}$, the
log-linear temperature is $\theta$ rather than $T$, the tariff interval set is $\mathcal{H}$
while the scalar horizon is $H$, and machine buffers keep the calligraphic
$\mathcal{B}_m$.}
\label{tab:notation}
\begin{tabular}{@{}ll@{}}
\toprule
Symbol & Meaning \\
\midrule
$\mathcal{J},\mathcal{M},\mathcal{V},\mathcal{L}$ & jobs, machines, vehicles, locations \\
$N=\mathcal{M}\cup\mathcal{V}$ & players of the cooperative game, $n=|N|$ \\
$O_{i,k}$, $p_{i,k}$ & operation $k$ of job $i$ and its processing time \\
$T_{i,k}$, $\tau_{i,k}$ & transport leg $k$ of job $i$ and its duration \\
$\mathcal{H}$, $c_h$ & tariff intervals and prices \\
$H$ & scalar horizon, and the hard deadline \\
$\lambda$ & horizon scaling factor (a feasibility parameter) \\
$B^{\min},B^{\mathrm{cap}},B^{\max}$ & SoC floor, ceiling, capacity (mAh) \\
$q^{\mathrm{idle}},q^{e},q^{\ell}$ & depletion rates, idle / empty / loaded (mAh/min) \\
$K_{CH}$ & charger capacity \\
$\Phi$, $\widehat{C}_{\max}$, $\widehat{E}$ & normalised objective and its two terms \\
$\omega$ & trade-off weight \\
$\mathcal{D}$ & finite ToU-aligned delay grid \\
$\alpha$, $a_p$, $a_p^{0}$ & joint profile, action of $p$, null action \\
$W$, $u_p$ & welfare (potential) and marginal-contribution utility \\
$\widehat{\Phi}(s_q,\alpha)$ & stage evaluation of the objective \\
$\pi_0$ & reference completion policy \\
$c(S)$, $v(S)$ & coalition cost and savings \\
$\varphi(v)$ & Shapley value \\
$\varepsilon^\star$ & least-core radius \\
$CS$, $S_{\max}$ & coalition structure and size cap \\
$\theta$ & log-linear learning temperature \\
\bottomrule
\end{tabular}
\end{table}

\section{Proofs}
\label{app:proofs}
Theorems~\ref{thm:t1} and~\ref{thm:t3} and Propositions~\ref{prop:n5}, \ref{prop:t6}
and~\ref{prop:scale} are proved where they are stated, since each proof is short and the
argument is easier to check next to the statement it establishes. This appendix records the
assumption lists.

\paragraph{Theorem~\ref{thm:t1} (exact stage potential).} Requires: (a) $\widehat{\Phi}$ is a
function of $(s_q,\alpha)$ alone --- guaranteed by fixing a deterministic execution rule and a
deterministic completion policy $\pi_0$; (b) finite action sets --- guaranteed by the finite
delay grid $\mathcal{D}$ and the finite task set; (c) \emph{sequential} deviations for the
finite-improvement argument. Does \emph{not} require submodularity, monotonicity or any
structure on $\Phi$: the identity is a cancellation.

\paragraph{Theorem~\ref{thm:t3} (price of stability).} Requires (a) only. It is a
\emph{stage-game} statement: $\mathrm{PoS}=1$ for each stage game, not for the trajectory.

\paragraph{Proposition~\ref{prop:n5} (congestion).} Requires a shared resource of finite
capacity that saturates, and a cheap tariff window whose capacity is exhausted by fewer
coordinating agents than are available. The instance in the proof satisfies both by
construction; the released artefact contains it as a regression test.

\paragraph{Proposition~\ref{prop:p5} (conditional convexity).} Requires the congestion-free
condition \emph{and} an unproved min-envelope step. Stated conditionally throughout; nothing
in the paper depends on it.

\paragraph{Proposition~\ref{prop:t6} (monotonicity).} Holds for the exact characteristic
function by a replication argument. It is not claimed for the single-sweep surrogate, whose
violation rate is measured instead.

\paragraph{Theorem~\ref{thm:leastcore} (least core).} Feasibility: the equal split satisfies
the equality constraint, and $\varepsilon$ can be driven low enough to satisfy every
inequality, so the programme is feasible. Boundedness: summing the constraints for
$S=N\setminus\{i\}$ over $i$ and using the equality bounds $\varepsilon$ above. The
equivalence $\varepsilon^\star\ge0 \Leftrightarrow \mathcal{C}(v)\neq\emptyset$ is immediate
from the definition of the core as the $\varepsilon=0$ feasible set.

\section{Benchmark data}
\label{app:bench}
The ten job sets and four travel-time layouts of~\citet{bilge1995} are reproduced in the
released artefact in the canonical instance format, together with the regenerated tariff
profile of every instance, its normalisation anchors, and its certified feasibility margin.
The regeneration rule is: period lengths $\ell_k=\mathrm{round}(H f_k)$ for $k=1,\dots,5$ with
$f=(3,5,3,3,2)/24$ and round-half-away-from-zero, the sixth period absorbing
$[b_5,H)$. This reproduces the published boundaries of periods one to five on all
\jnum{bench.instances} instances and, unlike the published table, never terminates the last
period after $H$.

\bibliographystyle{elsarticle-harv}
\bibliography{cas-refs,added-refs}

\end{document}
